% SAME-FAMILY ROOT WORKING FRAGMENT, 25 October 2026.
% Current complete batch: D17 with ordinary rank one only.
% B17/D17 actual-I bodies remain to append. Unintegrated/uncompiled.
\subsection{Rank-one spectators beside the rank-seventeen classical factors}
\label{app:rank17classical18}
The ordinary rank-one case and the actual isolated-short case use
different information. We first treat the ordinary \(D_{17}\)
alternative without imposing a shape on its rank-one opposite.

\subsubsection{\(D_{17}\) with an ordinary nonshort \(A_1\)}
\label{app:D17ordinary18}
Use the actual unit \(D_{17}\) roots, centre fifteen and forks
sixteen/seventeen. The long duals are \(\Omega_i=\sqrt2P_i\),
\(i\le15\); the fork duals are
\((P_{16}\mp e_{17})/\sqrt2\). The last coordinate is signed.
For the original decomposition \(z=E(p-x-ye)\), ordinary rank
one gives \(y^2>.1\), and hence \(\norm x^2<.9\).
Every classical cap retains this old \(e\), with full foreign
cost \(Q=\beta^2<1\), either zero or at least \(1/4\).

The valid bulk parameter is \(M=11\), so helper coverage first
reaches twelve, with an actual strong eleven/twelve. If five were
weak, its actual point supplies \(x_1>.5965,S_6>2.01815\).
The positive reference through one, six and that strong endpoint
has energy
\[
 >.7600299945+(\sqrt{11}-\sqrt6)^2/5>.9104024765.
\]
Its slopes decrease; the endpoint is below its actual strong floor.
This contradicts \(\norm x^2<.9\). Five is strong before any
new tail or head argument.

Here are the fresh tail matrices, to specify precisely which local
arguments of Appendix~\ref{app:exact18layers} are reused. At
thirteen/fourteen deletion retains \(A_{12}+D_4\) and
\(A_{13}+D_3\). In the actual \(D_4\), long-fourteen contact
three has \((\Delta,\sigma)=(1/2-Q,g+1/2)\); either fork-three
has \((1/2-Q,g+1/4)\). In the actual \(D_3\), middle-three,
end-three, end-four and end-five have respectively
\[
 (1/2-Q,g+1/2),\quad(5/8-Q,g+1/4),\quad
 (1/4-Q,g+1/(2\sqrt2)),\quad
 (1-3\phi^2/8,g+\phi/4).
\]
The sufficient strong ratio is \(<T(6)<25/16\).
These are the same finite local matrices, not a whole-component
theorem with its rank changed. Left-only and isolated caps use
their long-root clearings at \(D_i=2i\): every isolated foreign
case is strong, the pure five/seven/at-least-eight/dotted points
have their \(.61\) loss bounds, and a positive foreign far-three
is strong. Own-zero and owns three/four/six/parallel have the
actual affine/null or reciprocal-shadow dispositions of that
inventory, retaining every foreign cost. Both fork choices remain.

For completeness the right lists leave only these nondirect rows.
At thirteen, a fork-three with \(Q=0\) and own five survives;
all positive foreign cases have ratio \(>2.1\), and finite at
least seven or dotted has ratio \(>1.15\sqrt2\).
At fourteen, middle-three is direct. End-four at own zero has
kernel \((\sqrt2,2,2,1,1)\) on the cap, contacted fork,
centre, other fork and old fourteen; all positive owns are direct.
The sealed end-five \(H_4\) white block has ratio \(>2.99\)
even at own zero. End-three has its own-zero \(E_8\), own-three
cycle, own-four star, own-six path and parallel null. Positive
foreign cost is direct; \(Q=0,g\ge1\) has ratio at least
\(\sqrt{5/2}\). Thus only pure own five and finite at least
seven remain there.

Full bridges use simple-three left endpoints only. The
\(D_4\)/middle-\(D_3\) bridge has residual \(1/(2i)\)
and numerator at least \(g+1/4+1/(2i)\). The end-three bridge
at fourteen has residual \(1/8+1/28=9/56\) and numerator
at least \(g+1/4+1/28\). Their own-zero principals are actual
all-three cycles, and every positive own is directly strong.
A foreign arm at the internal bridge cannot be finite. No
exceptional finite type has the required rank. The surviving weak
rows are therefore pure far-left helpers and the explicit right
exceptions just listed.

The right own-five exception at \(i=13,14\) uses its actual
left \(H_4\) projected point with correction
\(g(3n_i/2+n_{i-1}+n_{i-2}/2)\), \(g=\phi/2\).
With \(\alpha=\sqrt2(g+1/4)\),
\[
 A^2=i\alpha^2-\kappa,\quad
 \kappa=1/2\ (13),\ 5/8\ (14),\quad
 G>\sqrt{i-\kappa/\alpha^2}-.091,
 \quad R<i.
\]
The normalized displacement is flat through \(i-3\). The finite
point has positive time and all original wall signs, including the
two forks and old rank-one component. Positive squares give
\(G>3.48\) and \(G>3.61\). For thirteen choose actual strong
nine/ten; the local one-square formula and the one ordinary charge
give totals \(>1.0757,>1.0762\). Coverage thus reaches thirteen.

For every remaining finite \(.9<g<1\) at fourteen, average the
two actual \(A_2\) cores with corrections
\(g(4n_{14}/3+2n_{13}/3)\) and
\(g(4n_{14}/3+2n_{15}/3)\). Both satisfy every wall inequality.
Their correction cross \(8g^2/9\) is retained in the average.
With \(a=4g^2/3-1\), its data are
\[
 A^2=14\alpha^2-5/8,\quad R<14,\quad
 G=(A-\sqrt a)/\alpha>
 \sqrt{69/5}-2\sqrt6/15>3.388.
\]
The derivative numerator \((5/8)/A-(g/3+1)/\sqrt a\) is
negative on the full interval, proving the endpoint bound. At
actual strong ten/eleven the local payment at its individual
\(1/7\) coefficient, plus \(y^2>.1\), gives totals
\(>1.0295,>1.0294\). No dotted parameter is missing: it was
directly strong. Coverage reaches fourteen with only pure helpers.

Their nonadjacency extends by the actual two-cap calculation:
adjacent helper products exceed \(1-ab>-11/25>-1/2\), so
the native gap forces orthogonality. The retained cap Schur matrix
is \(I-t\mathbf1\mathbf1^T\), \(t=(i-1)/(2i)\), with
positive inverse. The unchanged prefix replacement gives weighted
loss \(<7/16+1/(4\sqrt{10})<.61\), contradicting two weak
floors. Retained right blocks remain positive; no larger-rank
prefix hypothesis is invoked.
There is therefore a strong thirteen/fourteen reference with
energy \(q>.88864\), first-five mean \(b<.361\) and individual
\(P_5\) coefficient \(>.146\). This gives only the safe
\(q_{13}\) floor; it does not assert strong fourteen or a
\(q_{14}\) floor.

Delete long one. The full retained \(D_{16}\) permits only
isolation or a three at its long endpoint two. The latter has
\(A^2=2g^2+2g+Q\), \(C_1=(2g+1)/\sqrt{1-2Q}\).
Own zero has zero time or its positive full-D null; own three
has its triangle. Every other own gives \(x_1>.643\).
Its legal first-five graft has gain
\((5/4)(.643-b)^2>.0994\), already contradictory.

For an isolated-within-D head,
\(c=Ap-g\Omega_1-\beta e\), \(A^2=2g^2+Q-1\).
The actual foreign shadow and original-farthest orthogonal-opposite
obstruction give \(A>\beta\), so \(2g^2>1\).
Owns zero/three/four are impossible; six and parallel have actual
future nulls. All remaining own and foreign cases have one of the
following actual pure-D pressures. A normalized displacement
\(Z\), supported in the first five coordinates with sum one,
has \(\ip xZ\ge H\), \(R=\norm Z^2\); weight \(\mu\)
costs \(\mu/5\) and gains
\(2\mu(H-b)-\mu^2(R-1/5)\) in this one reference.

Pure five uses the finite \(H_4\) correction
\(g(3n_1/2+n_2+n_3/2)\). It gives \(Z=P_4/4\), \(R=1/4\),
\(H^2=1/(4\phi)>2/13>.392^2\). Weight \(.4\) gives
gain \(>.0168\). Pure finite at least seven uses
\(c+g^{-1}n_1+2(g^{-1}-g)n_2\), and, for
\(u=g^{-2}\in(1,1.232)\),
\[
 Z=(1-u/2)(e_1+e_2)+(u-1)e_3,\quad
 H={\sqrt{2-u}-\sqrt{7u-3u^2-4}\over\sqrt2},\quad
 R=3-4u+3u^2/2.
\]
The complete interval seals for weight \(1/2\) are
\[
\begin{array}{c|c|c|c}
u& H\text{ lower}&R\text{ upper}&\text{gain lower}\\\hline
[1,1.1]&.6708-.1871&.5&.0477\\
[1.1,1.2]&.63245-.204125&.415&.013575\\
[1.2,1.232]&.6196-.2&.36&.0186.
\end{array}
\]
They follow by positive squaring and the quadratic radical maximum
\(1/24\). The native-seven cubic gives the upper \(1.232\)
uniformly for all finite labels at least seven.

For dotted or foreign-cleared \(0<u=(1-Q)/g^2<1\), the positive
core \(c+\beta e+gu n_1\) gives exactly the \(Z,H\)
pressure in the common dotted-head calculation of
Appendix~\ref{app:BD16typed18}. Its scalar proof uses only this
pressure, \(b<.361\) and the individual first-five capacity.
It therefore gives the same legal gain \(>.03624\) here at
weight \(\sqrt2/5\); no typed-defect hypothesis or conclusion
is used. All original walls are retained by the positive core.
At \(u=1\) is a nonzero feasible future null.
For \(Q>0,u>1\), the native alphabet forces own five and
foreign three only. Its actual finite \(H_3\) projection on
\(c,n_1,e\) gives \(Z=P_2/2\), \(R=1/2\),
\(H^2=1/(2\sqrt5)>2/9>.47^2\). Weight \(1/2\) gives
gain \(>.034\). Every foreign finite label and the full dotted
own interval is thus included.

Each price is at most \(.1<.146\). Every selected finite point
is an actual positive projection; every negative core has future
time and nonpositive product with all original walls. All head
cases give a legal gain greater than \(.0135\). Hence
\[
 \norm x^2+y^2>.88864+.0135+.1=1.00214>1.
\]
This closes precisely the ordinary \(D_{17}\) rank-one case.
Tail elimination contradictions have not been added to the final
reference, and no short-\(I\) energy or arithmeticity is assumed.
